% ============================================================================
% Fichier  : partie4/P04_C14_declaration_applicabilite.tex
% Titre    : La Déclaration d'Applicabilité et le Plan de Traitement
% Auteur   : MINKA MI NGUIDJOÏ Thierry Emmanuel
% ============================================================================
% Sources : ISO/IEC 27001:2022 (clause 6.1.3) ; ISO/IEC 27005. Fictif : Eséka.
% ============================================================================

\chapter{La Déclaration d'Applicabilité et le Plan de Traitement}
\label{chap:dda}
\chaptermark{La déclaration d'applicabilité}

\epigraph{%
    « La déclaration d'applicabilité est la carte d'identité du SMSI~:\\
    elle dit ce qu'on fait, ce qu'on ne fait pas, et pourquoi. »%
}{Principe du cours}

\niveauChapitre{Intermédiaire}{Le Chapitre~\ref{chap:traitement}.}

\begin{boxobjectifs}
\begin{itemize}
    \item Définir la déclaration d'applicabilité (\sigle{DdA}) et son rôle.
    \item Construire une \sigle{DdA} à partir des risques et de l'Annexe~A.
    \item Distinguer \sigle{DdA} et plan de traitement des risques.
    \item Justifier l'inclusion et l'exclusion d'une mesure.
    \item Faire acter le risque résiduel par la direction.
\end{itemize}
\end{boxobjectifs}

\bigskip

% ============================================================================
\section*{Question directrice : comment prouver que les mesures retenues sont les bonnes ?}
% ============================================================================

L'auditeur ne croit pas sur parole~: il demande à voir le lien entre risques et
mesures. La déclaration d'applicabilité est le document qui établit ce lien.
Ce chapitre la construit.

\medskip
\textbf{Plan du chapitre.}
\begin{enumerate}
    \item Rôle de la DdA (section~\ref{sec:c14-role}).
    \item Construire la DdA (section~\ref{sec:c14-construire}).
    \item DdA et plan de traitement (section~\ref{sec:c14-plan}).
    \item Inclure, exclure (section~\ref{sec:c14-justifier}).
    \item Acter le risque résiduel (section~\ref{sec:c14-acter}).
\end{enumerate}

\separateur

% ============================================================================
\section{Rôle de la DdA}
\label{sec:c14-role}
% ============================================================================

\begin{boxdef}
La \termecle{déclaration d'applicabilité} (\sigle{DdA})~: document exigé par la
clause~6.1.3 de \norme{ISO/IEC 27001} qui, pour chaque mesure de l'Annexe~A,
indique si elle est \emph{retenue}, sa \emph{justification} (inclusion ou
exclusion) et son \emph{état} de mise en \oe uvre.
\end{boxdef}

La \sigle{DdA} est le document que tout auditeur demande en premier~: elle
résume le \sigle{SMSI} en une page de synthèse et une grille. Elle fait le pont
entre la partie IV (risques) et la partie V (mesures).

% ============================================================================
\section{Construire la DdA}
\label{sec:c14-construire}
% ============================================================================

\begin{tableaucours}{L{3.4cm}L{1.6cm}L{3.2cm}L{3.0cm}}{%
    \tableauHeader{Mesure (Annexe A)} &
    \tableauHeader{Retenue} &
    \tableauHeader{Justification} &
    \tableauHeader{État}
}
Sauvegarde & Oui & Risque R1 (état civil) & En place \\
Contrôle d'accès & Oui & Risque R2 (taxes) & En place \\
Chiffrement & Oui & Risque R2 & En cours \\
Développement sécurisé & Non & Pas de développement interne & Exclue \\
\end{tableaucours}
\vspace{-0.8em}
\begin{center}{\small\itshape Tableau~14.1. Extrait de la DdA de la mairie d'Eséka}\end{center}

Chaque ligne relie une mesure à un risque (inclusion) ou explique pourquoi la
mesure ne s'applique pas (exclusion). Rien n'est laissé implicite.

% ============================================================================
\section{DdA et Plan de Traitement}
\label{sec:c14-plan}
% ============================================================================

\begin{boxdef}
Le \termecle{plan de traitement des risques}~: document qui, pour chaque risque,
précise la stratégie, les mesures, le responsable et l'échéance. La \sigle{DdA}
dit \emph{quelles} mesures~; le plan dit \emph{comment et quand} les mettre en
\oe uvre.
\end{boxdef}

Les deux documents sont complémentaires~: la \sigle{DdA} est la photographie
des mesures~; le plan est la feuille de route. Ensemble, ils constituent le
« Plan » du \sigle{PDCA}.

% ============================================================================
\section{Inclure, Exclure}
\label{sec:c14-justifier}
% ============================================================================

\begin{boxprocedure}
\textbf{Constatation.} La mairie exclut la mesure « développement sécurisé » au
motif qu'elle n'écrit aucun logiciel elle-même.

\medskip
\textbf{Analyse.} L'exclusion est légitime et justifiée~: la mesure ne
correspond à aucun risque de la mairie. Mais la justification doit être écrite~:
une exclusion non justifiée est un écart d'audit. On ne retire une mesure que
si le risque associé n'existe pas.
\end{boxprocedure}

\begin{boxmetier}{Auditeur SMSI}
\textbf{Lire la DdA à rebours.} L'auditeur part de la \sigle{DdA} et remonte~:
cette mesure retenue répond-elle à un risque réel~? cette mesure exclue l'est-
elle à bon droit~? La \sigle{DdA} est le fil qui relie tout le \sigle{SMSI}.
\end{boxmetier}

% ============================================================================
\section{Acter le Risque Résiduel}
\label{sec:c14-acter}
% ============================================================================

Une fois la \sigle{DdA} et le plan établis, il reste un risque résiduel. La
clause~6.1 exige que les \termecle{propriétaires des risques} approuvent le plan
de traitement et \termecle{acceptent} le risque résiduel. À la mairie, c'est le
maire qui signe cette acceptation.

\begin{boxjuris}
Faire acter le risque résiduel par la direction n'est pas une formalité~: c'est
l'acte par lequel la responsabilité est assumée au bon niveau. En cas
d'incident, cette trace montre que le risque avait été identifié, pesé et
décidé, et non ignoré.
\end{boxjuris}

\separateur

\begin{boxresume}
\textbf{Ce qu'il faut retenir de ce chapitre~:}
\begin{itemize}
    \item La \textbf{DdA} indique, pour chaque mesure, si elle est retenue, sa
          justification et son état.
    \item Elle se \textbf{construit} à partir des risques et de l'Annexe~A.
    \item La \textbf{DdA} dit quelles mesures~; le \textbf{plan de traitement}
          dit comment et quand.
    \item \textbf{Inclusion} et \textbf{exclusion} se justifient par écrit.
    \item La direction \textbf{accepte} formellement le risque résiduel.
\end{itemize}
\end{boxresume}

\begin{boxinfo}
\textbf{Points de contrôle de l'élève.} Je sais~: définir la \sigle{DdA}~; la
construire~; la distinguer du plan de traitement~; justifier une inclusion ou
une exclusion~; expliquer l'acceptation du risque résiduel.
\end{boxinfo}

\bigskip

\begin{boxexo}[title={\ding{45}\quad Exercice~14.1. Rôle de la DdA \niveaubase}]
Qu'indique la \sigle{DdA} pour chaque mesure de l'Annexe~A~?
\end{boxexo}

\begin{boxexo}[title={\ding{45}\quad Exercice~14.2. DdA ou plan \niveaubase}]
Laquelle dit « quelles mesures », lequel dit « comment et quand »~?
\end{boxexo}

\begin{boxexo}[title={\ding{45}\quad Exercice~14.3. Exclusion \niveaubase}]
Pourquoi la mairie peut-elle exclure la mesure « développement sécurisé »~?
\end{boxexo}

\begin{boxexo}[title={\ding{45}\quad Exercice~14.4. Une ligne \niveaumoyen}]
Rédigez une ligne de \sigle{DdA} pour la mesure « journalisation ».
\end{boxexo}

\begin{boxexo}[title={\ding{45}\quad Exercice~14.5. Exclusion non justifiée \niveaumoyen}]
Pourquoi une exclusion non justifiée est-elle un écart d'audit~?
\end{boxexo}

\begin{boxexo}[title={\ding{45}\quad Exercice~14.6. Lire à rebours \niveaumoyen}]
Comment un auditeur se sert-il de la \sigle{DdA} pour vérifier le \sigle{SMSI}~?
\end{boxexo}

\begin{boxexo}[title={\ding{45}\quad Exercice~14.7. Acceptation \niveaumoyen}]
Qui, à la mairie, accepte le risque résiduel, et pourquoi la trace importe-t-
elle~?
\end{boxexo}

\begin{boxexo}[title={\ding{45}\quad Exercice~14.8. Complémentarité \niveaumoyen}]
Montrez que DdA et plan de traitement sont complémentaires.
\end{boxexo}

\begin{boxexo}[title={\ding{45}\quad Exercice~14.9. DdA de la mairie \niveauexpert}]
Construisez un extrait de \sigle{DdA} de huit mesures pour la mairie d'Eséka.
\end{boxexo}

\begin{boxexo}[title={\ding{45}\quad Exercice~14.10. Plan complet \niveauexpert}]
Rédigez le plan de traitement des risques de la mairie (cinq risques, stratégie,
mesures, responsable, échéance, acceptation).
\end{boxexo}

\bigskip

\begin{boxinfo}
\textbf{Pour aller plus loin.}
\begin{itemize}
    \item ISO/IEC 27001, clause~6.1.3~\cite{iso27001}~; ISO/IEC 27005~\cite{iso27005}.
\end{itemize}
\end{boxinfo}

\begin{boxinfo}
\textbf{Transition.} Le plan décidé, place à la mise en \oe uvre~: la partie V
s'ouvre au Chapitre~\ref{chap:orga} sur les mesures organisationnelles et
humaines.
\end{boxinfo}
